\chapter{Technologies and Tools Used}
\label{chap:tecno-ferra}

\section{Introduction}
\label{chap3:sec:intro}

This chapter describes the concrete set of technologies and tools on
which the work reported in this document is built. The set was not
chosen in a single decision, but converged over the course of the
project through a mixture of constraints (what the production
\ac{API} already uses, what the laboratory's network allows, what the
proposal explicitly mandates), comparisons (between equivalent
alternatives) and pragmatic preference. The chapter is structured by
layer of the stack: \ac{API} contracts and serialization formats in
Section~\ref{sec:apis}, the Java/Jakarta ecosystem in
Section~\ref{sec:java}, persistence and \ac{SQL} databases in
Section~\ref{sec:sql}, and servers, tooling and network in
Section~\ref{sec:servers}.

\section{APIs, Serialization Formats and Specification}
\label{sec:apis}

\subsection{REST and the Multi-format Response}

All endpoints exposed by the applications discussed in this report
follow the \ac{REST} architectural style as described by
Fielding~\cite{fielding2000rest}: resources are identified by stable
\acp{URL}, manipulated through a small set of \ac{HTTP} verbs (with
\texttt{GET} dominating for read-only \ac{GNSS} metadata), and the
\ac{API} is stateless on the server side. There is no session
affinity, which is a property that the differential test harness
described in Chapter~\ref{chap:caso-epos} relies on heavily: the same
request can be replayed against two independent server processes and
the responses can be compared directly.

The production \ac{API} (\texttt{api\_v4.0}) is unusual in that it
serves the \emph{same} resource in several serialization formats,
selected by a path parameter at the end of the \ac{URL}. The supported
formats are \ac{JSON} (default), \ac{XML} (produced through Jackson's
XML data format module), \ac{CSV} (produced through a custom writer)
and, for the file metadata endpoints, an executable shell
\emph{script} that downloads the corresponding files. This multi-format
output was inherited from the legacy \ac{API}, many of its clients
are scientific scripts that consume \ac{CSV} or \ac{JSON} directly,
and was preserved in the refactor; it is also the reason the
\ac{HTTP}-level differential harness has to compare not only headers
and status codes but also the body, format by format.
Figure~\ref{fig:multiformat} sketches the content negotiation pattern
that turns a single resource method into four wire formats.

\begin{figure}[h]
  \centering
  \resizebox{\linewidth}{!}{%
  \begin{tikzpicture}[
      every node/.style={font=\small},
      box/.style={draw, rounded corners=2pt, align=center, thick,
                  minimum height=9mm, minimum width=22mm},
      writer/.style={box, fill=dkgreen!10, minimum width=20mm},
      arrow/.style={-{Stealth[length=2mm]}, thick}
    ]
    \node[box, fill=dkblue!10] (req) {\texttt{GET /api/stations/.../\{format\}}};
    \node[box, fill=mauve!10, below=8mm of req, minimum width=46mm]
      (resource) {\textbf{StationEndpoints}\\[2pt]\footnotesize same Java method, regardless of format};
    \node[box, below=8mm of resource, minimum width=46mm, fill=lgray]
      (negot) {Content negotiation \\\footnotesize (path parameter \texttt{\{format\}})};

    \node[writer, below=10mm of negot, xshift=-36mm] (json)   {\textbf{JSON} \\\footnotesize Jackson};
    \node[writer, right=4mm of json]                 (xml)    {\textbf{XML}  \\\footnotesize jackson-xml};
    \node[writer, right=4mm of xml]                  (csv)    {\textbf{CSV}  \\\footnotesize custom writer};
    \node[writer, right=4mm of csv]                  (script) {\textbf{Script}\\\footnotesize files only};

    \node[box, below=10mm of csv, fill=dark!10, minimum width=66mm]
      (resp) {HTTP response body};

    \draw[arrow] (req)       -- (resource);
    \draw[arrow] (resource)  -- (negot);
    \draw[arrow] (negot)     -- (json);
    \draw[arrow] (negot)     -- (xml);
    \draw[arrow] (negot)     -- (csv);
    \draw[arrow] (negot)     -- (script);
    \draw[arrow] (json)      -- (resp);
    \draw[arrow] (xml)       -- (resp);
    \draw[arrow] (csv)       -- (resp);
    \draw[arrow] (script)    -- (resp);
  \end{tikzpicture}%
  }
  \caption{Multi-format content negotiation in \texttt{api\_v4.0}. The
  same Java resource method is reused across four serializations,
  selected by a path parameter, and the differential harness must
  compare the body of each format separately.}
  \label{fig:multiformat}
\end{figure}

\subsection{\ac{JSON} and \ac{XML}}

\ac{JSON} is the default serialization and the format on which most
of the test fixtures of this project rely. The Jakarta \ac{REST}
implementation provided by Payara takes care of object serialization
through Jackson, with negligible per-request overhead. \ac{XML} is
provided by the \texttt{jackson-dataformat-xml} module, configured to
preserve the same Java models that produce \ac{JSON}. This dual
configuration is valuable because it removes a whole class of
\emph{schema drift} bugs: as long as the Java model is correct, both
\ac{JSON} and \ac{XML} outputs are consistent with each other.

\subsection{\ac{CSV}}

\ac{CSV} is the format favoured by historical clients of the \ac{API}
and is produced through a small set of writer classes that flatten
the entity tree into rows. The differential harness handles \ac{CSV}
exactly the same way it handles \ac{JSON} and \ac{XML}: the textual
body of the response is treated as an opaque blob that must match,
byte for byte, against the baseline.

\subsection{\ac{OpenAPI} and Swagger \ac{UI}}

The contract of the production \ac{API} is described in an
\ac{OpenAPI} 3.1.0~\cite{openapi} document
(\texttt{openapi.json}) that ships inside the \ac{WAR}. The same
document is served to a Swagger \ac{UI}~\cite{swaggerui} instance
hosted at \texttt{/swagger.html}, which provides an interactive,
in-browser playground for the endpoints. Two practical benefits arise
from this setup. First, the document is the single source of truth
shared with downstream consumers of the \ac{API}, in particular the
\ac{EPOS} portal. Second, the same document can be used to drive
contract tests: a future evolution of the test framework, briefly
discussed in Chapter~\ref{chap:conc-trab-futuro}, is to derive the
\ac{HTTP} differential test cases automatically from the \ac{OpenAPI}
specification instead of maintaining them in a separate \ac{JSON} file.

\section{Java and Jakarta EE}
\label{sec:java}

The applications target Java SE 21~\cite{javaSE21} (the production
\ac{API} and the second prototype) and Java SE 17 (the first
prototype, kept compatible with the \texttt{maven:\allowbreak 3.9.6-eclipse-\allowbreak temurin-21}
\ac{CI} image). The choice of Java was effectively imposed by the
existing code base; \texttt{record} types are used liberally for
request/response \acp{POJO} and the new \texttt{switch} expression
syntax replaces several long \texttt{if/else} chains from the legacy
\ac{API}.

\textbf{Jakarta EE 11}~\cite{jakartaee} provides the platform layer.
The renaming from Java EE in 2018 had already been absorbed by the
time this project started; the current code base targets the
\texttt{jakarta.*} package namespace. The specifications actively
used across the three applications are: Jakarta \ac{REST} 4.0
(\ac{JAX-RS})~\cite{jaxrs} for endpoints, Jakarta \ac{CDI} 4.1 for
injection, Jakarta Persistence 3.2 (\ac{JPA})~\cite{jpa} for
\ac{ORM}, Jakarta Bean Validation 3.1 for declarative constraints,
Jakarta JSON Binding 3.0 for serialization, and Eclipse
MicroProfile~\cite{microprofile} 2.1 \ac{JWT} (used by the second
prototype) plus Health 4.0 for probes. Spring Boot~\cite{springBoot}
was considered but rejected: the existing code base is already on
Jakarta EE on top of an existing Payara installation, and a parallel
migration would have brought no proportional benefit.

\section{\ac{SQL} and the Database Layer}
\label{sec:sql}

The persistence model of the production \ac{API} is heavily
relational: \ac{GNSS} stations, networks, agencies, countries, file
types, observation summaries and embargo windows are all linked
through foreign keys, and most endpoints execute several joins per
request. \ac{SQL}~\cite{postgresql} is therefore the natural query
language, and even where \ac{JPA}'s object-oriented abstractions
(Section~\ref{sec:jpa-hibernate}) are used, what reaches the database
is plain \ac{SQL}. The laboratory standardised on
\textbf{PostgreSQL}~\cite{postgresql} for the \ac{EPOS} production
database; the choice over \textbf{MySQL}~\cite{mysql} was driven
mainly by PostgreSQL's adherence to the \ac{SQL} standard (advanced
\texttt{JOIN} types, common table expressions, the
\texttt{numeric} type for high-precision station coordinates) and by
its rich set of index types (B-tree, BRIN, GiST). The PostGIS
extension is not in use today (Section~\ref{sec:epos-db}) but
remains a realistic future addition. MySQL was nevertheless used in
the docker-compose configuration of the second prototype
(Section~\ref{sec:prototypes}) to validate that the application-level
code is not tied to a specific dialect.

\subsection{\ac{JPA}, Hibernate and the Persistence Unit}
\label{sec:jpa-hibernate}

\ac{JPA}~\cite{jpa} is the Jakarta specification for object-relational
mapping in Java. The production \ac{API} uses
Hibernate~\cite{hibernate} as the \ac{JPA} provider, configured
through a \texttt{persistence.xml} file. The persistence unit is
named \texttt{MeuPU} and is declared as
\texttt{RESOURCE\_LOCAL}, which means that transactions are managed
explicitly by application code rather than by the container, a
deliberate choice that keeps the differential test harness in full
control of when transactions begin and end.

Database connections are pooled through \textbf{HikariCP}~\cite{hikaricp}
with a small, conservative pool (2--15 connections, 20-second connect
timeout, 30-second idle, 5-second validation), tuned for the
read-mostly workload of the \ac{API} and for the relatively low
concurrency observed in production.

\subsection{The \ac{EPOS} Database in Practice}
\label{sec:epos-db}

The production database lives on a PostgreSQL 16 instance reachable
from the laboratory network at \texttt{192.168.168.30:5432}, and the
same schema serves two fixtures: \texttt{DB\_3.0\_J}, a small but
representative production subset used by the correctness diffs, and
\texttt{DB\_3.0\_F}, the full dataset used for volume-oriented and
performance experiments. The schema contains forty-two tables and
around fifty foreign-key constraints, grouped into a core domain
(\texttt{station}, \texttt{network}, \texttt{agency},
\texttt{country}, \texttt{rinex\_file}, \texttt{highrate\_rinex\_file},
\texttt{file\_type}), the junction tables that materialise its
many-to-many relationships, a family of seven quality-control
summary tables, and a handful of auxiliary tables (embargoes,
documents, licenses, infrastructure). Station coordinates are
plain \texttt{numeric}; spatial logic is currently performed in Java
by the \texttt{Geometry} package rather than by PostGIS, which is
not installed (the only non-default extension is \texttt{plpgsql}).
The twenty-five existing indexes are mostly B-trees on foreign-key
columns, with \texttt{fillfactor=90} on the read-heavy quality-control
tables. The next planned optimisation, and the reason the larger
\texttt{F} fixture was provisioned, is the introduction of BRIN
indexes on the date columns of \texttt{rinex\_file}: \ac{RINEX} files
are written in approximate date order, so the on-disk layout already
correlates with the indexed column, which is the precondition for a
small and fast BRIN index. The differential framework introduced in
Chapter~\ref{chap:caso-epos} is used to validate that this
optimisation does not change observable behaviour against the
\texttt{F} fixture.

\section{Servers, Tooling and Network}
\label{sec:servers}

\textbf{Payara Server}~\cite{payara} is the Jakarta EE application
server adopted by the laboratory; the production \ac{API} is deployed
onto an existing Payara installation. Its companion product,
\textbf{Payara Micro}, is a single self-contained \ac{JAR} that
launches a deployable \ac{WAR} from the command line and is the
engine on which the \ac{HTTP}-level differential harness relies (the
harness boots one Payara Micro per \ac{WAR} on different ports and
replays the same requests against each). Apache Tomcat~\cite{apacheTomcat}
was considered for the prototypes but rejected because it does not
ship a full Jakarta EE 11 runtime out of the box.

All three applications are built with \textbf{Apache Maven}; the
production \texttt{pom.xml} configures, among others, the
\texttt{maven-war-plugin} (which produces the deployable
\texttt{ROOT.war}), \texttt{maven-failsafe-plugin} (integration tests
named \texttt{*IT.java} with \ac{JUnit} XML output) and
\texttt{maven-dependency-plugin} (which stages the Payara Micro
\ac{JAR} into \texttt{target/payara-micro/} so the \ac{HTTP} harness
can spawn it as a subprocess).

\textbf{GitLab}~\cite{gitlabci} provides both the source-control
front end and the \ac{CI}/\ac{CD} platform. The GitLab server is a
self-managed installation~\cite{gitlabSelfManaged} running on a
\ac{VM} inside the laboratory; the GitLab Runner~\cite{gitlabRunner}
is a separate process on another \ac{VM} (tagged
\texttt{ubi-runner-.31}) that picks up jobs and executes them inside
Docker containers. That separation matters because the runner is the
machine that needs network access to the laboratory's database, not
the control plane. \textbf{Git}~\cite{chacon2014pro} itself is used
non-trivially in the pipeline: \texttt{git worktree} is what
materialises each baseline build alongside the \texttt{HEAD} build on
the runner's filesystem, without disturbing the working copy
(Section~\ref{sec:pipeline}). All three repositories are hosted on
\texttt{gitlab.com}, the production \ac{API} under the
\texttt{segalubi} group, the two prototypes under the personal
namespace of the author.

\textbf{Docker}~\cite{docker} appears in two roles. \ac{CI} jobs run
inside the official
\texttt{maven:\allowbreak 3.9.6-eclipse-\allowbreak temurin-21} image,
which guarantees a clean and reproducible toolchain regardless of the
runner host. The two prototypes additionally ship a
\texttt{docker-compose.yaml}~\cite{composeSpec} for local development
(PostgreSQL plus, in the first prototype, a pair of Payara Server
containers for the dual-node Arquillian tests). The production
\ac{API} is intentionally \emph{not} containerised: the deployment
model is to copy the \ac{WAR} onto an existing Payara installation,
and an image would add a moving part with no benefit (this is
revisited in Chapter~\ref{chap:conc-trab-futuro}).

Finally, the production database is reachable only inside the
laboratory's network on \texttt{192.168.168.30:5432}, which is why
the runner is pinned by tag and why developers operating from outside
connect through the laboratory's \ac{VPN}, maintained by the SEGAL
staff. Alternatives such as WireGuard~\cite{wireguard} were
considered as lighter-weight options for one-off remote sessions but
were not adopted.

\section{Conclusion}
\label{chap3:sec:concs}

The technology stack on which this project is built is in large part
imposed by its environment: Jakarta EE on Payara, PostgreSQL on the
internal network, GitLab as the self-managed \ac{CI} platform. Where
there was room for choice, between Jakarta EE and Spring Boot,
between Payara Server and Apache Tomcat, between PostgreSQL and
MySQL, between JUnit 5 and TestNG, the decisions documented in
this chapter were guided by a preference for fewer moving parts, by
the alignment with what the laboratory already runs in production
and by the desire to keep the differential testing harness as simple
as possible. The next chapter describes the iterative methodology by
which these technologies were combined into a working solution.
